\chapter{Dynamic Geodesic and Operator Learners, and the Filtering Target Contract}
\label{ch:bf-neural-ot-dynamic-target-contract}

The previous chapter treated direct map-learning families such as ICNN/Brenier and
Monge-gap approaches.  The present chapter changes the learned object again.
Instead of asking the network to learn one static transport map, it asks what
happens when the learned object is a dynamic transport path, a velocity field, or
an operator acting on families of measures.

This chapter is intentionally secondary.  Its job is not to become another core
implementation chapter.  Its job is to explain why these broader learned families
carry an even heavier target-status burden than retained-teacher acceleration or
static direct-map learning.

\section{Reader orientation: what changes here?}
\label{sec:bf-neural-ot-dynamic-orientation}

The easiest way to read this chapter is to keep the following contrast in mind.

\begin{itemize}
  \item The retained-teacher neural chapter learned how to solve a declared
  transport teacher faster.
  \item The direct-map chapter learned one final transport map.
  \item This chapter asks what happens when the learned object is richer than one
  final map: a path, a velocity field, or an operator on entire families of
  measures.
\end{itemize}

So the main burden is not just implementation difficulty.  It is target drift:
what scalar is now being computed, and what target would a later optimizer or HMC
routine actually be differentiating?

\section{Object inventory for dynamic and operator learners}
\label{sec:bf-neural-ot-dynamic-object}

A dynamic transport learner no longer tries to produce one final map
\(T\).  Instead it may try to learn:
\begin{itemize}
  \item a geodesic path \(\rho_\lambda\) between source and target measures,
  \item a velocity field \(v_\lambda(x)\) generating that path,
  \item or an operator \(\mathcal G\) that maps input measures to output
  trajectories or output measures.
\end{itemize}
This is a larger algorithmic change because the learned object is not merely a
replacement map.  It is a representation of how transport unfolds.

\paragraph{Reader checkpoint.}
The present chapter has crossed another object boundary.  A dynamic learner is not
just another static map family.  It tries to represent motion or evolution, not
only the endpoint transformation.

\section{Static map versus dynamic path: a toy contrast}
\label{sec:bf-neural-ot-dynamic-toy-contrast}

A tiny one-dimensional toy example makes the distinction concrete.  Suppose the
source cloud is
\[
  X_{\mathrm{src}}=\{-1,1\}
\]
and the target cloud is
\[
  X_{\mathrm{tgt}}=\{0,2\}.
\]
A static-map learner tries to learn one rule such as
\begin{equation}
\label{eq:bf-neural-ot-dynamic-static-map}
  T(x)=x+1,
\end{equation}
which moves the source points directly to the target points.

A dynamic path learner instead tries to describe how the particles move as a
function of pseudo-time \(\lambda\in[0,1]\).  One simple path is
\begin{equation}
\label{eq:bf-neural-ot-dynamic-path-example}
  x(\lambda)=x+\lambda,
  \qquad
  \lambda\in[0,1].
\end{equation}
At \(\lambda=0\) the source points are unchanged, and at \(\lambda=1\) they have
reached the target points.  The endpoint transformation is the same as the static
map, but the learned object is richer: it describes the whole path, not just the
final location.

This toy distinction is enough for the present chapter.  It shows why a dynamic
learner carries more structure, more opportunity, and more target risk than a
single direct map.

\section{Dynamic geodesic learners}
\label{sec:bf-neural-ot-dynamic-geodesic}

Some neural OT papers learn Wasserstein geodesics or associated velocity fields.
These are mathematically appealing because particle-flow filtering already has a
pseudo-time flavor.  The attraction is therefore obvious: perhaps one could learn
reusable transport dynamics instead of solving a fresh static problem each time.

The caution is equally important.  A dynamic geodesic learner does not answer the
same question as a retained finite Sinkhorn solver.  It learns a richer object and
must therefore be judged against a richer target contract.

\paragraph{Reader checkpoint.}
A geodesic learner is not just a faster static map and not just a warm-started
solver.  It is a learner for the entire transport path or for the dynamics
generating that path.  That greater ambition is why the chapter is more about
burden and verification than about a single near-implementable core algorithm.

\begin{algorithm}[ht]
\caption{\texttt{train\_dynamic\_transport\_path}}
\begin{algorithmic}[1]
\Require source-cloud sampler, target-cloud sampler, pseudo-time grid \(0=\lambda_0<\cdots<\lambda_K=1\), learned velocity-field or path parameterization, rollout integrator, endpoint loss, path regularizers, optimizer, stopping rule, and validation suite
\Ensure trained path or velocity-field parameters, deployable terminal map or path object, and pathwise validation diagnostics
\State Initialize the path or velocity parameters and optimizer state
\While{the stopping rule has not been met}
  \State Draw source and target cloud minibatches
  \State Initialize the source particles at pseudo-time \(\lambda_0\)
  \For{each pseudo-time step \(k=0,\ldots,K-1\)}
    \State Evaluate the learned velocity field or path increment at \(\lambda_k\)
    \State Advance the particles using the declared rollout or integration rule
  \EndFor
  \State Compute the endpoint loss against the target cloud at \(\lambda=1\)
  \State Compute any declared pathwise regularizers, such as smoothness or action penalties
  \State Backpropagate the total loss and update the parameters
\EndWhile
\State On held-out clouds, evaluate endpoint fit, path regularity, discretization sensitivity, and downstream scalar sensitivity
\State Return the trained path/velocity object together with the declared diagnostics
\end{algorithmic}
\end{algorithm}

\section{Operator-learning families}
\label{sec:bf-neural-ot-dynamic-operator}

A related family learns operators on spaces of measures.  In such approaches, the
input is itself a measure or empirical cloud, and the output is either another
measure, a path of measures, or a set of latent coordinates describing the
transport.  This is useful as a representation idea, but it moves even farther
from the finite retained-teacher transport layer of the previous chapters.

The main conceptual advantage is reuse across families of problems.  The main
conceptual cost is that the learned object is even less tied to a single declared
teacher solve.  That makes these methods valuable as research-direction chapters,
but not as the default first implementation route in this monograph.

\section{Adjacent boundary-marker families}
\label{sec:bf-neural-ot-dynamic-boundary-families}

Several nearby families remain useful only as boundary markers in this monograph:
- measure-to-measure regression operators,
- local Wasserstein regression,
- manifold or Riemannian neural OT.

These show that neural transport ideas extend beyond the current finite-cloud
Euclidean BayesFilter setting.  They do not, by themselves, answer the present
implementation question.

\section{Same-scalar rule revisited}
\label{sec:bf-neural-ot-dynamic-same-scalar}

The most important target-contract warning survives every change of learned object.
If a downstream routine uses a scalar \(\widetilde L\) built from a learned map,
path, or operator, then the reported gradient must differentiate that same scalar.
Even if this value-gradient parity holds, the result is coherent only for the
learned target unless a correction returns to the teacher or original target.

The small toy example from the earlier chapter already showed how a learned output
cloud can change a simple downstream scalar.  The same logic becomes more severe
here because the learned object is richer and therefore has more ways to drift from
the teacher semantics.  A path learner can change not only the final cloud, but
also the whole trajectory through which intermediate summaries are defined.

\paragraph{Plain-language takeaway.}
A smooth learned path is not automatically a faithful path for the original target.
It may simply be a coherent path for a different learned target.

\section{Consistency notes for richer learned objects}
\label{sec:bf-neural-ot-dynamic-promotion}

For richer learned objects, a few mathematical consistency checks are especially
important:
\begin{enumerate}[label=(\roman*)]
  \item whether endpoint or pathwise improvement is accompanied by stability of the
  downstream scalar,
  \item whether the learned object behaves robustly under weight, dimension, or
  particle-count shift,
  \item whether a reported gradient is the gradient of the same scalar value that
  is later interpreted downstream,
  \item whether any claim of target preservation is supported by an explicit
  correction mechanism rather than by smoothness or speed alone.
\end{enumerate}
These are not governance gates.  They are simply the mathematical checks needed to
understand what object is being computed and whether it remains identifiable from
the declared scalar and learned path or operator.

\begin{algorithm}[ht]
\caption{\texttt{audit\_dynamic\_transport\_target\_contract}}
\begin{algorithmic}[1]
\Require learned output object (map, path, or operator), downstream scalar-construction routine, reported gradient routine, optional teacher or corrected baseline scalar, shift-test suite, and declared tolerances for scalar and gradient mismatch
\Ensure a consistency report containing same-scalar parity status, shift-robustness status, any failed conditions, and the strongest mathematically supported interpretation of the learned object
\For{each evaluation problem in the declared suite}
  \State Construct the learned output object for that problem
  \State Compute the downstream scalar \(\widetilde L\) from that learned object
  \State Compute the reported gradient of that same scalar
  \If{a teacher or corrected baseline exists}
    \State Compute the baseline scalar and, where applicable, the baseline gradient
    \State Record scalar and gradient gaps relative to the baseline
  \EndIf
  \State Check that the reported gradient is tied to the same scalar \(\widetilde L\)
  \State Evaluate shift robustness under particle-count, weight, and context variation
  \State Record any failed scalar, gradient, or robustness condition
\EndFor
\State Return the full consistency report and the strongest interpretation of the learned object consistent with those results
\end{algorithmic}
\end{algorithm}

\paragraph{Reader checkpoint.}
This procedure does not approve or reject a method in an institutional sense.  It
simply organizes the mathematical question the chapter cares about: is the learned
object still identifiable through the scalar and gradient that later code actually
uses?

\section{Relation to the full HMC target chapter}
\label{sec:bf-neural-ot-dynamic-handoff}

This chapter stops at the target-status burden created by broader dynamic and
operator learners.  It does not settle the HMC question.  That full taxonomy still
belongs to Chapter~\ref{ch:bf-dpf-hmc-target-suitability}, which interprets the
scalar and gradient contracts rung by rung across the whole DPF ladder.

\section{What this chapter does and does not claim}
\label{sec:bf-neural-ot-dynamic-boundary}

This chapter explains why the broadest neural OT families should be read as
secondary research directions in this monograph.  It does \emph{not} claim:
\begin{itemize}
  \item that geodesic, operator, and adjacent measure/manifold families solve the
  same implementation problem as retained-teacher acceleration,
  \item that a richer learned object is automatically more useful for BayesFilter,
  \item that these families are already ready for the same implementation status as
  the earlier chapters.
\end{itemize}
Its role is to close the expanded OT block responsibly by keeping the broadest
learned families in view while preserving the main algorithmic center of gravity
in the earlier, more concrete chapters.

\FloatBarrier
